% % % % % % % % % % % % % % % % % % % % % % % % % % % % % % % % % % % % % % % % % % 
% EALS LATEX TEMPLATE
% ABSTRACT
% English version
% 
% % % % % % % % % % % % % % % % % % % % % % % % % % % % % % % % % % % % % % % % % % 
% % % Authors % % % 
% Jana Winkler (jana.winkler@student.tugraz.at)
%
% % % Changelog % % % 
% 0.1 - (10.11.2023) - added contents line

% % % To-Do % % % 
%
% % % % % % % % % % % % % % % % % % % % % % % % % % % % % % % % % % % % % % % % % % %

\chapter*{\centering Abstract}
Advances in processor hardware for voltage source inverters in recent years have enabled the use of more complex control schemes for field-oriented control (FOC) of permanent magnet synchronous machines (PMSM). Due to its elegant formulation for multiple input, multiple output systems and flexibility regarding extensions, model predictive control (MPC) is seen as an attractive approach. In contrast to the widely used FOC cascade of linear compensators as controllers with feedforward decoupling for nonlinearities, MPC allows for a direct consideration of the whole model, while the constraints imposed by the inverter and the motor can be directly incorporated through an optimization-based approach. Furthermore, smooth transitions between constrained and unconstrained operating regions can be ensured without the need for anti-windup measures.\\
This work aims to design a continuous control set model-predictive current controller for an interior magnet PMSM driven by a two-stage voltage-source inverter, which can be gradually integrated into existing inverter control systems. The mathematical model of the PMSM exhibits several nonlinearities that complicate the use of various control approaches. In order to meet the strict real-time constraints imposed by the implementation while still being able to use an accurate prediction model, the MPC optimization problem, which is at the core of this control approach, is derived in detail, and the various assumptions made in the process are explained.\\
Another key focus of this thesis is the rejection of disturbances, such as those caused by model faults. To this aim, approaches based on an embedded integrator, weighting of integral errors in the cost function, and disturbance estimation are compared; ultimately, an application-specific formulation of the first approach is used, which does not require an incremental prediction model for the control of electric drives. Furthermore, the typical MPC extensions for this application, namely, dead-time compensation due to computation and switching dead times, and dynamic constraints on the stator current, are derived.\\
After the selected controller implementation has been discussed and validated via simulation, an experimental evaluation of the controller is conducted in the machine laboratory at Graz University of Technology.

\chapter*{\centering Kurzfassung}
Durch den Fortschritt in Prozessorhardware für Frequenzumrichter in den letzten Jahren wurde es möglich, komplexere Regelungssysteme für die feldorientierte Regelung (FOC) von permanenterregten Synchronmaschinen (PMSM) zu verwenden. Aufgrund ihrer eleganten Formulierung bezüglich Mehrgrößensystemen und Flexibilität bezüglich Erweiterungen ist die modellprädiktive Regelung (MPC) ein attraktiver Ansatz. Im Gegensatz zur weit verbreiteten FOC Kaskade von linearen  Gliedern, mit Entkopplung für nichtlineare Systemteile, als Regler, kann bei MPC das gesamte Streckenmodell mitsamt  Beschränkungen durch Umrichter und Maschine optimierungsbasiert direkt berücksichtigt werden. Weiters  können ohne Anti-Windup Maßnahmen glatte Übergänge zwischen beschränkten und unbeschränkten Betriebsbereich gewährleistet werden.\\
Diese Arbeit richtet sich dem Ziel, einen modellprädikitven Stromregler für die Innenmagnet PMSM, gespeist durch einen Zweipunkt-Spannungszwischenkreisumrichter mit kontinuierlicher Modulationsstufe zu entwerfen, welcher schrittweise in bestehende Umrichtersteuerungen übernommen werden kann. Das mathematische Modell der PMSM zeigt mehrere Nichtlinearitäten, welche die Verwendung Regleransätze erschweren. Um die implementierungsbedingten harten Echtzeitbedingungen einhalten zu können und trotzdem ein akkurates Prädiktionsmodell verwenden zu können wird das  Optimierungsproblems des MPC, welches der Kern dieses Regleransatzes ist, detailliert hergeleitet und diverse dabei getroffene Annahmen begründet.\\
Ein weiterer Kernpunkt dieser Arbeit widmet sich der Unterdrückung von Störungen, beispielsweise aufgrund von Modellfehlern. Dazu werden Ansätze basierend auf einem eingebetten Integrator, Gewichtung der Integralfehler in der Kostenfunktion und Störungsschätzung miteinander verglichen und schlussendlich eine anwendungsspezifische Formulierung des ersten Ansatzes verwendet, der für die Regelung von elektrischen Antrieben ohne inkrementelles Prädiktionsmodell auskommt.  Weiters werden die typischen MPC Erweiterungen für besagte Anwendung, nämlich Totzeitkompensation aufgrund der Berechnungs--und Schalttotzeit und dynamische Beschränkungen des Statorstroms hergeleitet.\\
Nachdem die gewählte Reglerimplementierung diskutiert und per Simulation validiert wurde, folgt die experimentelle Evaluierung dieser in dem Maschinenlabor der Technischen Universität Graz.
